\section{Related Work}
\label{sec:related_work}

\paragraph{Reinforcement Learning with Verifiable Rewards.}
Reinforcement learning from human feedback~\citep{ouyang2022training} aligns language models with human preferences through learned reward models, while reinforcement learning with verifiable rewards (RLVR) replaces preference models with programmatic correctness checks for tasks whose final answers are objectively verifiable~\citep{shao2024deepseekmath,deepseekr1,lambert2024tulu3}.
Recent RLVR work improves GRPO-style optimization through dynamic sampling and clipping~\citep{yu2025dapo}, extended training schedules~\citep{liu2025prorl}, finer credit assignment~\citep{kazemnejad2025vineppo}, and asymmetric handling of negative samples~\citep{zhu2025negsample}. These methods mainly operate on outcome-level reward signals, where structurally distinct traces can remain reward-equivalent. \method stays within the RLVR setting but injects deterministic topology and continuity diagnostics from latent reasoning DAGs, targeting structural credit assignment directly rather than through learned reward surrogates.

\paragraph{Process Reward Models for Reasoning.}
Process reward models (PRMs) score intermediate reasoning steps rather than only final answers, with early variants relying on dense human step annotations~\citep{lightman2023lets}.
Later work reduces annotation cost with Monte Carlo or tree-search rollouts~\citep{wang2024math,luo2024omegaprm,zhang2025qwen25mathprm}, while generative PRMs model verifier reasoning explicitly~\citep{thinkprm2025,genprm2026,she2025rprm}. Implicit PRM~\citep{yuan2024implicitprm} derives step-level signals from outcome feedback, yet PRMBench~\citep{song2025prmbench} shows persistent weaknesses on fine-grained error categories. \method is complementary to this line: it does not train a neural verifier, and it does not claim proof-level verification; instead, it introduces deterministic global-structure supervision over dependency topology and local continuity.

\paragraph{On-Policy Distillation and Reasoning Compression.}
Knowledge distillation~\citep{hinton2015distilling} compresses large models into smaller students, and on-policy distillation (OPD) closes the train-test gap by letting the student generate its own trajectories under teacher supervision~\citep{minillm2023,agarwal2024gkd,song2026opdsurvey}.
Recent OPD variants for reasoning include OPSDC~\citep{opsdc2026}, ExOPD~\citep{yang2026exopd}, BOND~\citep{sessa2024bond}, and KDRL~\citep{xu2025kdrl}; black-box alternatives such as GAD~\citep{ye2025gad} further broaden applicability. However, most methods still optimize token-level matching objectives without explicit structural control over dependency preservation. In parallel, interpretability studies show graph-like latent organization in long-CoT reasoning~\citep{besta2024graph,zhong2026chains2dags}. \method bridges these lines by using one deterministic DAG diagnostic interface for both reward optimization and on-policy distillation, so compression is conditioned on structural adequacy rather than only token imitation.
